\clearpage\thispagestyle{empty}
\pdfbookmark[0]{Quellen und Versübersetzungen zur Mollweide-Besprechung}{mollweide1814-vi-sources}
\begin{center}\Large\bfseries Quellen und Versübersetzungen zur Mollweide-Besprechung\end{center}

Die KI-Übertragung des deutschen Originals in heutiges Deutsch folgt der Besprechung von Carl Friedrich Gauß in den \emph{Werken}, Bd.~VI (Göttingen, 1874), S.~567--569. Der Erstdruck erschien in den \emph{Göttingischen gelehrten Anzeigen}, Stück~71 vom 2.~Mai 1814, S.~705--710. Im Sammelband endet die Besprechung in Band~VI nach der ersten Abhandlung. Für die weiteren Abhandlungen und die Briefe verweist sie auf Band~IV (1873), S.~361--362. Die Formen Fomahand und Fohamond folgen dem Sammelband.

Karl Brandan Mollweides \emph{Commentationes mathematico-philologicae tres} erschienen 1813 in Leipzig bei Carl Cnobloch. Das unten verlinkte Exemplar befindet sich in der ETH-Bibliothek Zürich (Rar~21664). Die erste Abhandlung zitiert die Stelle aus den \emph{Georgica} IV auf S.~4 und behandelt auf S.~28--30 andere Deutungen, darunter die von Samuel Horsley.

\noindent\textbf{Der in der Besprechung zitierte Titel.} Drei mathematisch-philologische Abhandlungen zur Erklärung zweier schwieriger Stellen, der einen bei Vergil und der anderen bei Platon, sowie zur Prüfung zweier Messregeln Columellas. Beigefügt ist ein Brief des Verfassers dieser Abhandlungen, Karl Brandan Mollweide, Professor der Astronomie an der Universität Leipzig, an den hochangesehenen J.~G. Schneider über die geometrischen Auszüge des Epaphroditus und des Vitruvius Rufus.

\noindent\textbf{Vergil, \emph{Georgica} IV.231--235.}
\begin{quote}\small Zweimal bringen sie reichen Ertrag ein. Es gibt zwei Erntezeiten,\\
sobald die Plejade Taygete der Erde ihr schönes Antlitz zeigt\\
und die verschmähten Fluten des Ozeans mit dem Fuß zurückstößt\\
oder wenn sie, vor dem Gestirn des wasserreichen Fisches fliehend,\\
trauriger vom Himmel in die winterlichen Wogen hinabsteigt.\end{quote}

\noindent\textbf{Vergil, \emph{Georgica} I.215--218.}
\begin{quote}\small\rightskip=-2em Im Frühling sät man Bohnen. Dann nehmen die lockeren Furchen\\
auch dich auf, Luzerne, und die jährliche Pflege der Hirse beginnt,\\
wenn der leuchtende Stier mit seinen vergoldeten Hörnern das Jahr eröffnet\\
und der Hund untergeht, dem gegenüberstehenden Gestirn \emph{weichend}.\end{quote}

Erstdruck (\href{https://images.p.sub.uni-goettingen.de/iiif/image/gdz:PPN319721507_1814:00000709/full/1650,/0/default.jpg}{705}, \href{https://images.p.sub.uni-goettingen.de/iiif/image/gdz:PPN319721507_1814:00000710/full/1650,/0/default.jpg}{706}, \href{https://images.p.sub.uni-goettingen.de/iiif/image/gdz:PPN319721507_1814:00000711/full/1650,/0/default.jpg}{707}, \href{https://images.p.sub.uni-goettingen.de/iiif/image/gdz:PPN319721507_1814:00000712/full/1650,/0/default.jpg}{708}).

\href{https://www.e-rara.ch/zut/content/titleinfo/15849974}{Katalogeintrag}, \href{https://www.e-rara.ch/i3f/v20/15849980/full/1400,/0/default.jpg}{Titelblatt}, \href{https://www.e-rara.ch/i3f/v20/15849989/full/1400,/0/default.jpg}{Vergil bei Mollweide, S.~4}. Mollweides Ergebnisse und andere Deutungen, S.~\href{https://www.e-rara.ch/i3f/v20/15850013/full/1400,/0/default.jpg}{28}, \href{https://www.e-rara.ch/i3f/v20/15850014/full/1400,/0/default.jpg}{29}, \href{https://www.e-rara.ch/i3f/v20/15850015/full/1400,/0/default.jpg}{30}.

\clearpage
